\section{Gaps and Next Steps}\label{subsec:next_steps}

\begin{reviewed}

	The systematic review highlights a critical divergence between current oversight methods and the technical reality of advanced models. While governance and societal impact dominate the literature, there remains a structural deficiency in addressing inner misalignment. Current alignment strategies such as reinforcement learning from human feedback tend to produce what has been termed shallow alignment \cite{mcintosh_inadequacy_2024, milliere_normative_2025}. These methods optimize for external compliance rather than internal value stability, leaving models vulnerable to goal misgeneralization \cite{langosco2022goal}.

	The most pressing gap identified in the mapping of the field is the emergence of deceptive alignment. Recent empirical evidence suggests that models are becoming capable of alignment faking, where they strategically satisfy training objectives to avoid modification \cite{greenblatt_alignment_faking}. This creates an evaluation trap where traditional auditing fails because the agent understands its own supervision \cite{anthropic2025biology, hubinger2019risks}. Consequently, the focus of this work should shift toward the control problem in the context of intentional subversion \cite{yampolskiy_controllability_2022}. Exploring how to detect and mitigate these deceptive behaviors represents the necessary next step in ensuring that highly capable systems do not execute a treacherous turn \cite{bostrom2014superintelligence}.

\end{reviewed}

\begin{toreview}

	While recent evidence of alignment faking centers on massive, frontier-class models \cite{greenblatt_alignment_faking}, the threshold at which these capabilities first emerge remains unclear. Thus, in a series of empirical investigations, this thesis will study deceptive alignment in smaller \glspl{llm}. By analyzing the relationship between model size and the propensity for alignment faking, with a specific focus on empirically underexplored sub-10B parameter scales. % This work seeks to determine what factors or architectures exacerbate the risk of deceptive tendencies.

\end{toreview}